\documentclass[11pt]{article}

\usepackage[margin=1in]{geometry}
\usepackage{amsmath,amssymb,amsfonts,mathtools}
\usepackage{booktabs}
\usepackage{longtable}
\usepackage{array}
\usepackage{enumitem}
\usepackage{hyperref}
\usepackage{xcolor}
\usepackage{titlesec}

\hypersetup{
    colorlinks=true,
    linkcolor=blue,
    urlcolor=blue,
    citecolor=blue
}

\setlist[itemize]{leftmargin=1.5em}
\setlist[enumerate]{leftmargin=1.8em}

\title{\textbf{Course Breakdown: Time Series Analysis}\\
\large EN.553.439 / EN.553.639 --- Spring 2026}
\author{Based on the course syllabus and required text:\\
J. D. Cryer and K.-S. Chan, \textit{Time Series Analysis with Applications in R}, 2nd ed.}
\date{}

\begin{document}

\maketitle

\section{Course Overview}

\subsection{Course Identity}

\begin{itemize}
    \item \textbf{Course title:} Time Series Analysis
    \item \textbf{Course numbers:} EN.553.439 / EN.553.639
    \item \textbf{Credits:} 3 credits
    \item \textbf{Primary perspective:} Time-domain modeling and frequency-domain analysis
    \item \textbf{Primary computational environment:} R
    \item \textbf{Required textbook:} Cryer and Chan, \textit{Time Series Analysis with Applications in R}, 2nd edition
\end{itemize}

\subsection{Course Description}

This course develops statistical methods for analyzing observations ordered in time.
The central challenge is that time-series observations are generally dependent:
an observation today may be related to observations from previous time points.
The course therefore extends familiar statistical ideas---such as means,
variances, regression, estimation, and prediction---to settings with serial
dependence, trends, seasonality, nonstationarity, and changing variability.

The course covers descriptive analysis, autocorrelation, stationary stochastic
processes, linear time-series models, ARMA and ARIMA models, model
identification, parameter estimation, diagnostic checking, forecasting,
seasonal models, regression with correlated errors, intervention analysis,
outlier detection, ARCH/GARCH volatility models, and selected spectral or
nonlinear time-series topics.

\section{Prerequisite Structure}

\subsection{Formal Prerequisites}

The formal prerequisites require prior coursework in probability/statistics and
multivariable calculus. The syllabus specifies the following combinations:

\begin{itemize}
    \item At least one statistics or probability course:
    \[
    \text{EN.553.310, EN.553.311, EN.553.420, EN.553.421, or EN.553.620}
    \]
    \item At least one calculus course:
    \[
    \text{AS.110.201, AS.110.212, or EN.553.291}
    \]
\end{itemize}

\subsection{Required Background Knowledge}

Before beginning the course, students should be comfortable with the following
topics.

\begin{longtable}{p{0.24\textwidth} p{0.34\textwidth} p{0.34\textwidth}}
\toprule
\textbf{Area} & \textbf{Prerequisite concepts} & \textbf{Why it is needed} \\
\midrule
\endhead

Probability &
Random variables, probability distributions, expectation, variance, covariance,
conditional expectation, joint distributions, independence &
Time-series models describe dependent random variables indexed by time. \\

Mathematical statistics &
Estimators, bias, consistency, likelihood, maximum likelihood estimation,
hypothesis tests, confidence intervals, asymptotic normality &
Used for ARMA/ARIMA parameter estimation, residual testing, and forecast
uncertainty. \\

Linear algebra &
Vectors, matrices, matrix multiplication, inverses, linear systems,
eigenvalues, eigenvectors, positive-definite covariance matrices &
Useful for regression, covariance calculations, multivariate formulations,
and spectral methods. \\

Calculus &
Differentiation, partial derivatives, optimization, logarithms, exponential
functions, Taylor approximations &
Needed for likelihood maximization, transformations, and asymptotic arguments. \\

Regression &
Simple and multiple linear regression, least squares, residuals, interpretation
of coefficients, model assumptions &
Required for deterministic trend estimation, regression with autocorrelated
errors, intervention analysis, and pre-whitening. \\

Programming in R &
Data import, vectors, matrices, plotting, functions, basic simulation,
model-fitting workflows &
The required text uses R examples, and computation is essential for fitting,
checking, and forecasting time-series models. \\

\bottomrule
\end{longtable}

\subsection{Recommended Pre-Course Review}

Students should review the following before the first two weeks:

\begin{enumerate}
    \item Expectation, variance, covariance, and correlation.
    \item Conditional probability and conditional expectation.
    \item Normal distributions and multivariate normal distributions.
    \item Ordinary least squares and residual analysis.
    \item Basic likelihood-based inference.
    \item Matrix notation for regression:
    \[
    \mathbf{y} = X\boldsymbol{\beta} + \boldsymbol{\varepsilon}.
    \]
    \item Basic R workflows for data visualization and regression.
\end{enumerate}

\section{Global Learning Objectives}

By the end of the course, students should be able to:

\begin{enumerate}
    \item Describe and visualize a time series using plots, summary statistics,
    transformations, and decomposition ideas.

    \item Define, estimate, and interpret the mean function, variance function,
    autocovariance function, autocorrelation function (ACF), and partial
    autocorrelation function (PACF).

    \item Distinguish deterministic trends from stochastic trends and choose
    appropriate methods for removing or modeling nonstationary behavior.

    \item Explain and assess strict stationarity and weak stationarity.

    \item Formulate, analyze, and interpret linear processes, autoregressive
    models, moving-average models, and ARMA models.

    \item Use differencing and transformations to construct suitable ARIMA
    models for nonstationary series.

    \item Identify plausible time-series models using the ACF, PACF, EACF,
    stationarity conditions, invertibility conditions, and domain knowledge.

    \item Estimate model parameters using least squares, conditional least
    squares, Yule--Walker equations, and maximum likelihood methods.

    \item Perform residual diagnostics and assess whether fitted models leave
    substantial serial dependence unexplained.

    \item Produce point forecasts and prediction intervals from ARIMA and
    seasonal ARIMA models.

    \item Build and interpret seasonal Box--Jenkins models.

    \item Analyze regression problems with serial correlation, interventions,
    outliers, and potentially spurious associations.

    \item Explain conditional heteroskedasticity and fit or interpret ARCH/GARCH
    volatility models.

    \item Recognize the connection between time-domain and frequency-domain
    approaches to time-series analysis.

\end{enumerate}

\section{Major Course Structure}

\subsection{Recommended Module Sequence}

The course is organized into thirteen modules. Modules 1--10 correspond most
closely to the core progression of the required textbook. Modules 11--13 cover
advanced applied material and topics identified in the syllabus.

\begin{longtable}{p{0.08\textwidth} p{0.25\textwidth} p{0.28\textwidth} p{0.29\textwidth}}
\toprule
\textbf{Module} & \textbf{Title} & \textbf{Major topics} & \textbf{Primary outcome} \\
\midrule
\endhead

1 & Foundations and Description &
Time ordering, visualization, mean, variance, covariance, ACF &
Describe dependence and structure in observed data. \\

2 & Stationarity and Basic Models &
Strict stationarity, weak stationarity, white noise, dependence &
Determine whether a stochastic process has stable probabilistic properties. \\

3 & Trend and Smoothing &
Deterministic trend, stochastic trend, regression trend estimation, smoothing &
Separate systematic long-run movement from short-run dependence. \\

4 & Linear Processes and ARMA &
Linear filters, AR, MA, ARMA, causality, invertibility &
Represent and interpret stationary serial dependence using linear models. \\

5 & Nonstationarity and ARIMA &
Differencing, integration, random walks, transformations, ARIMA &
Convert nonstationary series into modelable stationary components. \\

6 & Model Identification &
ACF, PACF, EACF, parsimonious model selection &
Propose plausible ARIMA-family structures from empirical evidence. \\

7 & Parameter Estimation &
Yule--Walker, least squares, conditional least squares, MLE &
Estimate model parameters and quantify uncertainty. \\

8 & Diagnostics and Validation &
Residual ACF, whiteness tests, model adequacy, fit assessment &
Determine whether the fitted model is adequate for inference and forecasting. \\

9 & Forecasting &
Linear prediction, ARIMA forecasting, forecast errors, intervals &
Generate and communicate point forecasts and prediction intervals. \\

10 & Seasonal Time Series &
Seasonal differencing, SARIMA, seasonal diagnostics, forecasting &
Model recurring periodic dependence and seasonal structure. \\

11 & Regression and Interventions &
Regression with correlated errors, intervention analysis, outliers,
spurious correlation, pre-whitening &
Analyze explanatory variables and structural changes in dependent data. \\

12 & Conditional Heteroskedasticity &
ARCH, GARCH, volatility clustering, conditional variance forecasting &
Model time-varying uncertainty and volatility. \\

13 & Extensions &
Nonlinear time series, spectral analysis, frequency-domain methods &
Recognize situations where linear time-domain models are insufficient. \\

\bottomrule
\end{longtable}

\section{Detailed Module Breakdown}

\subsection{Module 1: Time-Series Foundations and Descriptive Tools}

\subsubsection{Core Concepts}

\begin{itemize}
    \item Time series as an ordered sequence:
    \[
    \{Y_t : t \in \mathcal{T}\}.
    \]
    \item Observed series versus stochastic process.
    \item Time plots, lag plots, seasonal plots, and distributional summaries.
    \item Mean function:
    \[
    \mu_t = \mathbb{E}(Y_t).
    \]
    \item Variance function:
    \[
    \operatorname{Var}(Y_t).
    \]
    \item Autocovariance function:
    \[
    \gamma(h) = \operatorname{Cov}(Y_t,Y_{t+h}).
    \]
    \item Autocorrelation function:
    \[
    \rho(h) = \frac{\gamma(h)}{\gamma(0)}.
    \]
    \item Sample ACF and its interpretation.
\end{itemize}

\subsubsection{Learning Objectives}

Students will be able to:

\begin{itemize}
    \item Explain why independence assumptions often fail for sequential data.
    \item Construct and interpret time plots and autocorrelation plots.
    \item Calculate sample mean, variance, covariance, and autocorrelation.
    \item Distinguish short-lag, long-lag, negative, positive, and seasonal
    dependence patterns.
    \item Identify visible trends, changing variance, outliers, and periodicity.
\end{itemize}

\subsubsection{Difficult Topics}

\begin{itemize}
    \item Understanding that autocorrelation is indexed by \emph{lag}, not by
    two unrelated variables.
    \item Distinguishing correlation caused by a shared trend from genuine
    short-run serial dependence.
    \item Interpreting sample ACF values as estimates with sampling variability.
\end{itemize}

\subsubsection{Prerequisite Connections}

Probability, covariance, correlation, exploratory data analysis, and basic R
plotting.

\subsection{Module 2: Stationarity and Basic Stochastic Models}

\subsubsection{Core Concepts}

\begin{itemize}
    \item Strict stationarity.
    \item Weak or covariance stationarity.
    \item White noise and Gaussian white noise.
    \item Constant mean and constant variance.
    \item Autocovariance depending only on lag:
    \[
    \operatorname{Cov}(Y_t,Y_{t+h}) = \gamma(h).
    \]
    \item Ergodicity as a conceptual connection between time averages and
    probabilistic expectations.
\end{itemize}

\subsubsection{Learning Objectives}

Students will be able to:

\begin{itemize}
    \item State the definitions of strict and weak stationarity.
    \item Explain why weak stationarity is central to classical ARMA analysis.
    \item Determine whether simple processes are stationary.
    \item Identify white-noise behavior from model definitions and diagnostics.
    \item Explain why nonstationary series require transformation or specialized
    models.
\end{itemize}

\subsubsection{Difficult Topics}

\begin{itemize}
    \item The distinction between strict stationarity and weak stationarity.
    \item Understanding that a series can exhibit random fluctuations while
    retaining stable distributional properties.
    \item Recognizing that stationarity does not imply independence.
\end{itemize}

\subsubsection{Prerequisite Connections}

Joint distributions, moments, covariance, Gaussian distributions, and
independence.

\subsection{Module 3: Trends, Regression, and Smoothing}

\subsubsection{Core Concepts}

\begin{itemize}
    \item Deterministic trend models:
    \[
    Y_t = m(t) + X_t,
    \]
    where \(m(t)\) is a deterministic function of time.
    \item Linear, polynomial, and nonlinear deterministic trends.
    \item Stochastic trends.
    \item Trend estimation by regression.
    \item Smoothing and moving averages.
    \item Trend removal and residual analysis.
    \item Consequences of fitting stationary models directly to trending data.
\end{itemize}

\subsubsection{Learning Objectives}

Students will be able to:

\begin{itemize}
    \item Distinguish deterministic and stochastic trends.
    \item Fit and interpret a deterministic trend model.
    \item Assess whether detrending produces approximately stationary residuals.
    \item Explain when smoothing is useful and when it can conceal important
    temporal structure.
    \item Compare trend regression with differencing as approaches to
    nonstationarity.
\end{itemize}

\subsubsection{Difficult Topics}

\begin{itemize}
    \item Choosing between a deterministic trend model and a differenced model.
    \item Avoiding overfitting with high-order polynomial trends.
    \item Understanding why regression residuals may still be autocorrelated.
\end{itemize}

\subsubsection{Prerequisite Connections}

Linear regression, least squares, residuals, model diagnostics, and
polynomial functions.

\subsection{Module 4: Linear Processes and ARMA Models}

\subsubsection{Core Concepts}

\begin{itemize}
    \item Backshift operator:
    \[
    B Y_t = Y_{t-1}.
    \]
    \item General linear processes:
    \[
    Y_t = \sum_{j=0}^{\infty}\psi_j Z_{t-j}.
    \]
    \item Autoregressive processes:
    \[
    Y_t = \phi_1Y_{t-1}+\cdots+\phi_pY_{t-p}+Z_t.
    \]
    \item Moving-average processes:
    \[
    Y_t = Z_t+\theta_1Z_{t-1}+\cdots+\theta_qZ_{t-q}.
    \]
    \item ARMA\((p,q)\) processes:
    \[
    \phi(B)Y_t = \theta(B)Z_t.
    \]
    \item Causality, stationarity, invertibility, and roots of characteristic
    polynomials.
    \item ACF and PACF signatures of AR and MA models.
\end{itemize}

\subsubsection{Learning Objectives}

Students will be able to:

\begin{itemize}
    \item Express AR, MA, and ARMA models in both ordinary and backshift
    notation.
    \item Interpret AR and MA coefficients in terms of persistence and shock
    propagation.
    \item Determine stationarity for simple AR models.
    \item Explain invertibility and its importance for unique model
    representation.
    \item Derive or recognize ACF/PACF patterns associated with AR and MA
    processes.
\end{itemize}

\subsubsection{Difficult Topics}

\begin{itemize}
    \item Solving characteristic equations and applying root conditions.
    \item Distinguishing causality from invertibility.
    \item Understanding why AR ACFs typically decay while MA ACFs cut off.
    \item Translating between recursive and infinite linear-process
    representations.
\end{itemize}

\subsubsection{Prerequisite Connections}

Infinite series, polynomial algebra, recursion, covariance calculations, and
linear filters.

\subsection{Module 5: Nonstationarity, Transformations, and ARIMA Models}

\subsubsection{Core Concepts}

\begin{itemize}
    \item First differencing:
    \[
    \nabla Y_t = (1-B)Y_t = Y_t-Y_{t-1}.
    \]
    \item Higher-order differencing:
    \[
    \nabla^dY_t=(1-B)^dY_t.
    \]
    \item Random walk and random walk with drift.
    \item Integrated processes.
    \item ARIMA\((p,d,q)\) models:
    \[
    \phi(B)(1-B)^dY_t=\theta(B)Z_t.
    \]
    \item Variance-stabilizing transformations.
    \item Logarithmic and power transformations.
    \item Over-differencing and under-differencing.
\end{itemize}

\subsubsection{Learning Objectives}

Students will be able to:

\begin{itemize}
    \item Diagnose trend-based nonstationarity.
    \item Apply differencing appropriately to stabilize a series' mean behavior.
    \item Explain the role of the integration order \(d\) in ARIMA models.
    \item Select and interpret transformations for changing variance.
    \item Distinguish a trend-stationary process from a difference-stationary
    process.
\end{itemize}

\subsubsection{Difficult Topics}

\begin{itemize}
    \item The conceptual distinction between deterministic and stochastic trends.
    \item Interpreting unit-root-like behavior and random walks.
    \item Avoiding unnecessary differencing, which can introduce artificial
    negative autocorrelation.
    \item Correctly translating forecasts on a transformed or differenced scale
    back to the original scale.
\end{itemize}

\subsubsection{Prerequisite Connections}

Stationarity, linear models, transformations, regression trends, and
polynomial algebra.

\subsection{Module 6: Model Specification and Identification}

\subsubsection{Core Concepts}

\begin{itemize}
    \item Box--Jenkins iterative modeling philosophy.
    \item ACF-based identification.
    \item Partial autocorrelation function (PACF).
    \item Extended autocorrelation function (EACF).
    \item Parsimony.
    \item Candidate ARIMA model selection.
    \item Information criteria as supplementary model-comparison tools.
    \item Incorporating domain knowledge and interpretability.
\end{itemize}

\subsubsection{Learning Objectives}

Students will be able to:

\begin{itemize}
    \item Use time plots, ACFs, and PACFs to identify candidate ARIMA structures.
    \item Distinguish approximate cutoff and tail-off patterns.
    \item Explain why identification is probabilistic rather than mechanical.
    \item Propose multiple plausible models when empirical evidence is
    ambiguous.
    \item Choose parsimonious candidate models for estimation and validation.
\end{itemize}

\subsubsection{Difficult Topics}

\begin{itemize}
    \item Distinguishing population ACF/PACF patterns from noisy sample plots.
    \item Identifying mixed ARMA models, whose signatures are less clean than
    pure AR or pure MA models.
    \item Integrating differencing decisions with ARMA order selection.
    \item Avoiding selection based solely on a single plot or a single
    criterion.
\end{itemize}

\subsubsection{Prerequisite Connections}

ACF, PACF, ARMA models, stationarity, differencing, and exploratory plotting.

\subsection{Module 7: Estimation of Time-Series Models}

\subsubsection{Core Concepts}

\begin{itemize}
    \item Yule--Walker equations for autoregressive models.
    \item Least squares estimation.
    \item Conditional least squares estimation.
    \item Maximum likelihood estimation.
    \item Conditional versus exact likelihood.
    \item Standard errors and uncertainty quantification.
    \item Parameter constraints related to stationarity and invertibility.
\end{itemize}

\subsubsection{Learning Objectives}

Students will be able to:

\begin{itemize}
    \item Set up Yule--Walker equations for autoregressive models.
    \item Compare least-squares and likelihood-based estimation approaches.
    \item Interpret estimated ARMA/ARIMA coefficients and standard errors.
    \item Explain the role of initial conditions in conditional estimation.
    \item Recognize when optimization or numerical methods are needed.
\end{itemize}

\subsubsection{Difficult Topics}

\begin{itemize}
    \item Deriving and solving Yule--Walker equations.
    \item Understanding the distinction between conditional and exact likelihood.
    \item Likelihood optimization with correlated observations.
    \item Connecting parameter estimates to stationarity and invertibility
    constraints.
\end{itemize}

\subsubsection{Prerequisite Connections}

Matrix algebra, regression estimation, likelihood theory, numerical
optimization, and asymptotic inference.

\subsection{Module 8: Diagnostic Checking and Model Adequacy}

\subsubsection{Core Concepts}

\begin{itemize}
    \item Residuals and innovations.
    \item Residual time plots.
    \item Residual ACF and residual PACF.
    \item Tests for residual autocorrelation.
    \item Model adequacy and remaining structure.
    \item Distributional checks and outlier assessment.
    \item Iterative model refinement.
\end{itemize}

\subsubsection{Learning Objectives}

Students will be able to:

\begin{itemize}
    \item Compute and visualize residuals from a fitted model.
    \item Determine whether residuals resemble white noise.
    \item Identify evidence of remaining serial correlation, seasonality,
    heteroskedasticity, or nonlinearity.
    \item Explain why good in-sample fit does not guarantee a useful forecasting
    model.
    \item Revise a model based on diagnostic evidence.
\end{itemize}

\subsubsection{Difficult Topics}

\begin{itemize}
    \item Interpreting residual correlation tests while accounting for estimated
    parameters.
    \item Distinguishing random residual variation from meaningful model
    failure.
    \item Reconciling conflicting diagnostics, such as good residual ACF but
    poor forecasts.
\end{itemize}

\subsubsection{Prerequisite Connections}

Hypothesis testing, ACF, ARIMA estimation, residual analysis, and model
selection.

\subsection{Module 9: Forecasting and Prediction Intervals}

\subsubsection{Core Concepts}

\begin{itemize}
    \item One-step-ahead and multi-step-ahead forecasts.
    \item Linear prediction.
    \item Forecasts from AR, MA, ARMA, and ARIMA models.
    \item Forecast errors.
    \item Forecast error variance.
    \item Prediction intervals.
    \item Forecast evaluation.
    \item Recursive forecasting.
\end{itemize}

\subsubsection{Learning Objectives}

Students will be able to:

\begin{itemize}
    \item Produce forecasts from fitted ARIMA-family models.
    \item Interpret a point forecast as a conditional expectation.
    \item Construct and interpret prediction intervals.
    \item Explain why forecast uncertainty generally increases with forecast
    horizon.
    \item Compare forecasts using forecast errors and practical relevance.
\end{itemize}

\subsubsection{Difficult Topics}

\begin{itemize}
    \item Multi-step forecasting for ARMA and ARIMA models.
    \item Computing or interpreting forecast error variance.
    \item Returning forecasts to the original scale after differencing or
    transformation.
    \item Distinguishing confidence intervals for parameters from prediction
    intervals for future observations.
\end{itemize}

\subsubsection{Prerequisite Connections}

Conditional expectation, ARIMA models, variance propagation, normal
approximations, and model diagnostics.

\subsection{Module 10: Seasonal Box--Jenkins Models}

\subsubsection{Core Concepts}

\begin{itemize}
    \item Seasonal period \(s\).
    \item Seasonal differencing:
    \[
    (1-B^s)Y_t.
    \]
    \item Seasonal autoregressive and moving-average operators.
    \item SARIMA notation:
    \[
    \text{ARIMA}(p,d,q)\times(P,D,Q)_s.
    \]
    \item Seasonal ACF and PACF behavior.
    \item Seasonal diagnostics and forecasting.
\end{itemize}

\subsubsection{Learning Objectives}

Students will be able to:

\begin{itemize}
    \item Identify recurring seasonal behavior in a time plot and ACF.
    \item Apply seasonal differencing when appropriate.
    \item Specify a SARIMA model using seasonal and nonseasonal components.
    \item Diagnose seasonal residual dependence.
    \item Generate forecasts that preserve both short-term dynamics and
    recurring seasonal patterns.
\end{itemize}

\subsubsection{Difficult Topics}

\begin{itemize}
    \item Distinguishing seasonal effects from long-run trend.
    \item Choosing between ordinary and seasonal differencing.
    \item Interpreting multiplicative SARIMA notation.
    \item Reading ACF/PACF patterns at seasonal lags such as \(s, 2s,\) and
    \(3s\).
\end{itemize}

\subsubsection{Prerequisite Connections}

ARIMA modeling, differencing, model identification, residual diagnostics, and
forecasting.

\subsection{Module 11: Regression, Interventions, Outliers, and Pre-Whitening}

\subsubsection{Core Concepts}

\begin{itemize}
    \item Regression models with temporally dependent errors.
    \item Intervention analysis.
    \item Pulse, step, and gradual intervention effects.
    \item Outlier detection.
    \item Additive outliers and temporary changes.
    \item Spurious correlation between trending series.
    \item Cross-correlation analysis.
    \item Pre-whitening.
\end{itemize}

\subsubsection{Learning Objectives}

Students will be able to:

\begin{itemize}
    \item Explain why ordinary regression can be misleading when observations
    are serially correlated.
    \item Model the effect of an event or policy change on a time series.
    \item Identify and assess possible outliers or structural disruptions.
    \item Recognize spurious regression caused by common trends.
    \item Use pre-whitening conceptually to study lagged relationships between
    two dependent series.
\end{itemize}

\subsubsection{Difficult Topics}

\begin{itemize}
    \item Separating genuine intervention effects from ordinary stochastic
    variation.
    \item Distinguishing causation, association, and shared trends.
    \item Understanding why naive cross-correlations can be deceptive.
    \item Applying pre-whitening correctly before interpreting cross-correlation
    functions.
\end{itemize}

\subsubsection{Prerequisite Connections}

Multiple regression, ARIMA models, residual analysis, cross-covariance,
hypothesis testing, and forecasting.

\subsection{Module 12: ARCH and GARCH Models}

\subsubsection{Core Concepts}

\begin{itemize}
    \item Conditional heteroskedasticity.
    \item Volatility clustering.
    \item ARCH models.
    \item GARCH models.
    \item Conditional variance:
    \[
    \sigma_t^2 = \alpha_0 + \sum_{i=1}^{q}\alpha_i a_{t-i}^2
    + \sum_{j=1}^{p}\beta_j\sigma_{t-j}^2.
    \]
    \item Conditional variance forecasting.
    \item Mean modeling versus variance modeling.
\end{itemize}

\subsubsection{Learning Objectives}

Students will be able to:

\begin{itemize}
    \item Recognize time-varying variance and volatility clustering in data.
    \item Explain the difference between unconditional and conditional variance.
    \item Specify and interpret basic ARCH and GARCH models.
    \item Explain how GARCH models differ from ARIMA models.
    \item Produce or interpret conditional volatility forecasts.
\end{itemize}

\subsubsection{Difficult Topics}

\begin{itemize}
    \item Understanding that serial dependence can occur in squared residuals
    even when raw residuals are uncorrelated.
    \item Interpreting parameter constraints needed for positive conditional
    variance and stationarity.
    \item Separating mean dynamics from volatility dynamics.
\end{itemize}

\subsubsection{Prerequisite Connections}

Conditional expectation, variance, likelihood estimation, residual diagnostics,
and ARMA-style recursion.

\subsection{Module 13: Nonlinear and Frequency-Domain Extensions}

\subsubsection{Core Concepts}

\begin{itemize}
    \item Limitations of linear time-domain models.
    \item Nonlinear time-series behavior.
    \item Cycles and periodic components.
    \item Frequency-domain representation.
    \item Spectral density.
    \item Periodogram.
    \item Relation between autocovariance and spectral density.
\end{itemize}

\subsubsection{Learning Objectives}

Students will be able to:

\begin{itemize}
    \item Explain why some series cannot be adequately modeled using linear
    ARIMA models.
    \item Interpret frequency as cycles per unit time.
    \item Describe the purpose of a periodogram and spectral density.
    \item Connect repeated periodic behavior in a time plot to frequency-domain
    analysis.
    \item Recognize when spectral methods may complement time-domain methods.
\end{itemize}

\subsubsection{Difficult Topics}

\begin{itemize}
    \item Moving between time-domain autocovariance concepts and
    frequency-domain representations.
    \item Understanding the periodogram as a noisy estimator.
    \item Distinguishing deterministic seasonality, stochastic cycles, and
    spectral peaks.
\end{itemize}

\subsubsection{Prerequisite Connections}

Trigonometric functions, complex-number familiarity, covariance functions,
linear algebra, and signal-processing intuition.

\section{Concept Dependency Map}

\subsection{Foundational Dependencies}

\begin{enumerate}
    \item Probability and covariance concepts support ACF, autocovariance,
    stationarity, and white-noise models.

    \item Stationarity supports ARMA theory, model identification, parameter
    estimation, and linear forecasting.

    \item Trend analysis and transformations support nonstationary modeling and
    ARIMA construction.

    \item ARMA models support ARIMA models.

    \item ACF, PACF, and ARIMA theory support model identification.

    \item Identification supports estimation.

    \item Estimation supports diagnostic checking.

    \item Diagnostic checking supports reliable forecasting.

    \item ARIMA forecasting supports seasonal forecasting and intervention
    models.

    \item Regression and ARIMA concepts support pre-whitening and
    intervention analysis.

    \item Residual analysis and likelihood methods support ARCH/GARCH modeling.

    \item Autocovariance and linear-process concepts support spectral analysis.
\end{enumerate}

\subsection{Recommended Learning Path}

\[
\text{Probability/Regression}
\rightarrow
\text{Descriptive Analysis}
\rightarrow
\text{Stationarity}
\rightarrow
\text{ARMA}
\rightarrow
\text{ARIMA}
\rightarrow
\text{Identification}
\rightarrow
\text{Estimation}
\rightarrow
\text{Diagnostics}
\rightarrow
\text{Forecasting}
\rightarrow
\text{Seasonal/Regression/Volatility Extensions}.
\]

\section{Most Challenging Topics}

The following topics are likely to require the greatest amount of practice,
because they combine theoretical definitions, algebraic manipulations,
statistical inference, and computational interpretation.

\begin{longtable}{p{0.25\textwidth} p{0.28\textwidth} p{0.37\textwidth}}
\toprule
\textbf{Difficult topic} & \textbf{Why it is challenging} & \textbf{Recommended preparation} \\
\midrule
\endhead

Stationarity and unit-root behavior &
Requires separating deterministic trends, stochastic trends, and stationary
fluctuations. &
Practice classifying simple processes and compare detrending with
differencing. \\

ARMA causality and invertibility &
Depends on backshift algebra, characteristic roots, and infinite-series
representations. &
Review polynomial roots, geometric series, and AR(1)/MA(1) examples. \\

ACF/PACF model identification &
Sample plots are noisy, and mixed ARMA models do not have perfectly clean
patterns. &
Fit multiple candidate models and validate using residual diagnostics rather
than relying only on visual patterns. \\

Yule--Walker and likelihood estimation &
Combines covariance equations, matrix algebra, optimization, and statistical
inference. &
Review matrix notation and likelihood estimation for normal models. \\

Forecasting ARIMA models &
Requires correctly reversing differencing, propagating uncertainty, and
interpreting prediction intervals. &
Work through one-step and multi-step forecasts by hand before relying on
software. \\

Seasonal ARIMA &
Contains both ordinary and seasonal operators, which makes specification and
interpretation more complex. &
Master ordinary ARIMA notation before introducing seasonal components. \\

Pre-whitening and spurious regression &
Requires understanding how serial dependence can create misleading
cross-correlations and regression relationships. &
Review regression assumptions, residual autocorrelation, and common-trend
problems. \\

ARCH/GARCH &
Models conditional variance rather than only conditional mean, often requiring
a new way to interpret residual behavior. &
Review conditional expectation, variance, squared residuals, and likelihood
estimation. \\

Spectral analysis &
Changes the viewpoint from lags in time to frequencies and periodic
components. &
Review trigonometry, periodic functions, and the relationship between
covariance and cycles. \\

\bottomrule
\end{longtable}

\section{Assessment-Aligned Course Plan}

\subsection{Assessment Structure}

The syllabus specifies the following course-grade components:

\begin{itemize}
    \item Six homework assignments: \(25\%\).
    \item Midterm 1: \(25\%\).
    \item Midterm 2: \(25\%\), noncumulative.
    \item Midterm 3: \(25\%\), noncumulative.
\end{itemize}

\subsection{Suggested Assessment Alignment}

\begin{longtable}{p{0.18\textwidth} p{0.36\textwidth} p{0.34\textwidth}}
\toprule
\textbf{Assessment} & \textbf{Likely module coverage} & \textbf{Skills emphasized} \\
\midrule
\endhead

Homework 1 &
Modules 1--2 &
Descriptive analysis, moments, autocovariance, ACF, stationarity. \\

Homework 2 &
Modules 3--4 &
Trend modeling, smoothing, AR/MA/ARMA definitions, stationarity and
invertibility. \\

Midterm 1 &
Modules 1--4 &
Foundations, stationarity, trends, linear processes, ARMA theory. \\

Homework 3 &
Modules 5--6 &
Differencing, transformations, ARIMA specification, ACF/PACF/EACF
identification. \\

Homework 4 &
Modules 7--8 &
Yule--Walker estimation, likelihood ideas, model fitting, residual
diagnostics. \\

Midterm 2 &
Modules 5--8 &
ARIMA modeling, identification, estimation, and diagnostic checking. \\

Homework 5 &
Modules 9--10 &
Forecasting, prediction intervals, seasonality, SARIMA models. \\

Homework 6 &
Modules 11--12 &
Regression with correlated errors, interventions, outliers, pre-whitening,
ARCH/GARCH. \\

Midterm 3 &
Modules 9--12, with possible Module 13 material &
Forecasting, seasonal models, regression extensions, volatility models, and
selected advanced topics. \\

\bottomrule
\end{longtable}

\section{Suggested Weekly Course Progression}

\begin{longtable}{p{0.10\textwidth} p{0.29\textwidth} p{0.48\textwidth}}
\toprule
\textbf{Week} & \textbf{Primary focus} & \textbf{Expected capabilities} \\
\midrule
\endhead

1 & Introduction and descriptive tools &
Visualize time series; compute and interpret basic summary measures. \\

2 & Autocovariance, ACF, and stationarity &
Identify serial dependence and explain weak stationarity. \\

3 & Trends and smoothing &
Fit deterministic trends; compare trend removal and smoothing. \\

4 & AR, MA, and ARMA models &
Represent stationary linear models and interpret ACF/PACF behavior. \\

5 & Review and Midterm 1 &
Demonstrate mastery of foundational concepts through ARMA modeling. \\

6 & Differencing, transformations, and ARIMA &
Build stationary representations of nonstationary data. \\

7 & Identification and estimation &
Use ACF/PACF/EACF; estimate candidate models. \\

8 & Diagnostics and validation &
Assess residuals and revise inadequate models. \\

9 & Review and Midterm 2 &
Integrate ARIMA specification, estimation, and diagnostic checking. \\

10 & Forecasting &
Produce point forecasts and prediction intervals. \\

11 & Seasonal ARIMA &
Identify and model seasonal patterns. \\

12 & Regression and intervention methods &
Analyze covariates, events, outliers, and cross-correlations. \\

13 & ARCH/GARCH and Midterm 3 &
Model volatility and demonstrate mastery of later-course material. \\

14 & Nonlinear and spectral extensions &
Connect time-domain analysis to frequency-domain and advanced methods. \\

\bottomrule
\end{longtable}

\section{Summary of Course Outcomes}

Upon successful completion of the course, a student should be able to carry out
an end-to-end time-series analysis workflow:

\begin{enumerate}
    \item Plot and describe the observed series.
    \item Assess trend, seasonality, variance stability, and serial dependence.
    \item Transform or difference the series if needed.
    \item Identify plausible ARIMA or SARIMA model structures.
    \item Estimate parameters using appropriate statistical methods.
    \item Check whether residuals resemble white noise.
    \item Revise the model if diagnostics reveal remaining structure.
    \item Produce forecasts and prediction intervals.
    \item Interpret results in the context of the underlying application.
    \item Recognize when extensions such as intervention analysis, GARCH, or
    spectral analysis are more appropriate than a basic ARIMA model.
\end{enumerate}

\end{document}